\lecture{Лекция 7}{}

\subsection{Матрица проектирования}

Рассматривается задача
\[
  \begin{cases}
    \varphi(x) \to \min , & x \in \R^n , \\
    Cx = d , & C \ \text{размера } l \times n,\ d \in \R^{l} , \\
    F(x) \le g , & F \ \text{размера } m \times n,\ g \in \R^{m} ,
  \end{cases}
\]
где $\varphi$ может быть нелинейной. К такой задаче применяется \textbf{метод
Розена}. Всюду далее $\varphi(x)$ — выпуклая и непрерывно дифференцируемая.

\begin{definition}\label{def:projector}
Матрица $P$ размера $n \times n$ называется \textbf{матрицей проектирования}
(проектором), если
\[
  1^\circ\ P = P\T ,
  \qquad\qquad
  2^\circ\ P \cdot P = P .
\]
\end{definition}

\begin{example}\label{ex:projector}
Пусть $A$ — матрица размера $m \times n$ полного строчного ранга. Тогда
\[
  P = \underbrace{A\T}_{n \times m}
      \underbrace{\big(A A\T\big)^{-1}}_{m \times m}
      \underbrace{A}_{m \times n}
  \qquad (n \times n)
\]
— матрица проектирования.
\end{example}

\begin{proof}
$1^\circ$ Транспонируем, пользуясь $(XYZ)\T = Z\T Y\T X\T$:
\[
  \Big(A\T \big(A A\T\big)^{-1} A\Big)\T
  = A\T \Big(\big(A A\T\big)^{-1}\Big)\T \big(A\T\big)\T
  = A\T \Big(\big(A A\T\big)\T\Big)^{-1} A
  = A\T \big(A A\T\big)^{-1} A = P .
\]
Здесь использовано, что матрица $A A\T$ симметрична,
$\big(A A\T\big)\T = \big(A\T\big)\T A\T = A A\T$, а обратная к симметричной сама
симметрична: $\big(X^{-1}\big)\T = \big(X\T\big)^{-1}$.

$2^\circ$
\[
  P \cdot P
  = A\T \big(A A\T\big)^{-1}
    \underbrace{A \cdot A\T \big(A A\T\big)^{-1}}_{=\ E} A
  = A\T \big(A A\T\big)^{-1} A = P .
  \qedhere
\]
\end{proof}

\begin{idea}
Обратимость $A A\T$ — это и есть требование полного строчного ранга $A$: строки
матрицы ограничений должны быть линейно независимы. Для метода Розена это
означает, что среди активных ограничений нет дублирующих друг друга — иначе
матрицу нужно предварительно чистить.

Геометрически $P = A\T(AA\T)^{-1}A$ проектирует на линейную оболочку строк $A$,
а дополнительный проектор $Q = E - P$ (см. лемму ниже) — на \emph{ядро} $A$,
то есть на множество направлений, вдоль которых $Ax$ не меняется. Именно $Q$
и понадобится в методе.
\end{idea}

\begin{lemma}\label{lem:projector}
Пусть $P$ — матрица проектирования размера $n \times n$. Тогда
\begin{enumerate}[label=$\arabic*^\circ$,left=0pt]
  \item $Q \equiv E - P$ — матрица проектирования;
  \item $x\T P x \ge 0$ \ $\forall x$ (то есть $P$ положительно полуопределена);
  \item множества
        \[
          L \equiv \{\, y \mid y = Px \,\} ,
          \qquad
          L' \equiv \{\, z \mid z = Qx \,\}
        \]
        — ортогональные подпространства, причём
        \[
          \forall x \ \ \exists!\ \langle y, z\rangle :
          \quad x = y + z , \quad y \in L , \quad z \in L' .
        \]
\end{enumerate}
\end{lemma}

\begin{proof}
\textbf{1.1)} $Q\T = E\T - P\T = E - P = Q$.

\medskip
\textbf{1.2)}
$Q \cdot Q = (E-P)(E-P) = E - P - P + P\cdot P = E - P = Q$.

\medskip
\textbf{2)} Пользуясь обоими свойствами проектора,
\[
  x\T P x = x\T P \cdot P x = x\T P\T P x = \norm{P x}^2 \ \ge\ 0 .
\]

\medskip
\textbf{3.1) Ортогональность.} Что $L$ и $L'$ — подпространства, очевидно
(образы линейных отображений). Пусть $y \in L$, $z \in L'$, то есть
$y = Px_1$, $z = Qx_2$. Тогда
\[
  y\T z = (P x_1)\T (Q x_2) = x_1\T P\T Q\, x_2
        = x_1\T P (E - P)\, x_2
        = x_1\T \big(P - P\cdot P\big) x_2 = 0 .
\]

\medskip
\textbf{3.2) Существование и единственность разложения.} Существование:
\[
  x = Px + (E - P)x = Px + Qx ,
\]
и первое слагаемое лежит в $L$, второе — в $L'$. Единственность: пусть
\[
  x = y + z , \quad y \in L,\ z \in L' ,
  \qquad
  x = y_1 + z_1 , \quad y_1 \in L,\ z_1 \in L' .
\]
Вычитая, получаем $0 = (y - y_1) + (z - z_1)$, где $y - y_1 \in L$,
$z - z_1 \in L'$. Два ортогональных вектора в сумме дают нуль только если оба
нулевые, значит $y = y_1$, $z = z_1$.
\end{proof}

\inputfig{figures/fig-projection-decomposition}

\begin{idea}
Пункт~2$^\circ$ — именно то, что понадобится для доказательства убывания:
\[
  \nabla\T\varphi\, P\, \nabla\varphi = \norm{P\nabla\varphi}^2 \ \ge\ 0 ,
\]
причём строго больше нуля, пока проекция градиента ненулевая.

Пункт~3$^\circ$ — геометрический смысл всей конструкции: пространство
распадается в прямую ортогональную сумму $\R^n = L \oplus L'$, и проектор $P$
просто отбрасывает составляющую вдоль $L'$. В методе Розена $L'$ — это ядро
матрицы активных ограничений, то есть множество направлений, остающихся на
грани допустимого множества.
\end{idea}
